% ID: MAT-GAN-ELL-004
% Titolo: Corda e tangente a una ellisse
% Disciplina: Matematica
% Area: Geometria analitica
% Argomento: Ellisse
% Classe: Terza o quarta secondaria di secondo grado
% Difficolta: 4
% Tipo: problema_numerico
% Tag: matematica, geometria_analitica, ellisse
% Fonte: verifica_ellisse_anonimizzata
% Provenienza: materiale fornito dal docente; testo anonimizzato; soluzione redatta per Exergo

\begin{esercizio}
Dato il punto $P(1;\sqrt{3})$, sull'ellisse di equazione
    $3x^2+y^2=6,$
    determina:
    \begin{enumerate}
        \item la lunghezza della corda staccata dall'ellisse sulla retta passante per l'origine e per $P$;
        \item l'equazione della tangente all'ellisse in $P$.
    \end{enumerate}
\end{esercizio}

\begin{soluzione}
La retta OP è $y=\sqrt3x$. Sostituendo: $6x^2=6$, quindi gli estremi della corda
sono $P=(1,\sqrt3)$ e $Q=(-1,-\sqrt3)$; $PQ=\sqrt{4+12}=4$.
La tangente in P è $3x+\sqrt3y=6$.
\end{soluzione}
